\section{Consequences and scope}
\label{sec:discussion}

The construction gives a finite answer to the extraction problem for every
fixed positive exponent allowance. Its relative schemes are specified by
coordinate maps and integer parameters, and their recursive use includes
all finite overheads. Unrolling one fixed scheme and restricting its
coefficients to the rationals gives the seed in
Theorem~\ref{thm:main}. The argument uses the finite convolution relations
from \citet{OAI26}, but does not assume that preprint's spectral conclusion.
The theorem therefore provides a separate finite derivation of
\(\omega\le9/4\).

The quantitative family also resolves a denominator problem in exact
certificate constructions. Clearing the binary certificate's prescribed
rational denominators requires an exponentially growing scale. That
requirement belongs to those exact weights. Conserved rounding uses any
positive integer scale, preserves every tensor-factor balance, and loses at most
\(4J_a=O(a^3)\) in the entropy expression. The resulting family has
\(\log A=O(a^5(\log a)^2)\), with explicit constants and a proved choice
of \(a\) for every allowance. The binary denominator formula remains a useful comparison, but it is
not a size barrier for this construction or for matrix multiplication
algorithms in general.

The least-profile theorem identifies the exact limit of the stated scalar
axioms. The diagonal lower bound follows from tripling at the row index,
which is the subfamily used by the tensor maps; off-diagonal minimality
uses shifted tripling beyond the row index. Since the displayed profile itself
satisfies the axioms, those axioms alone cannot force faster diagonal
growth. The multiway relations improve the finite accounting while keeping
the same limiting profile exponent. A smaller asymptotic exponent would
require additional mathematical input.

Several components can also be used separately. The determinant map turns
a filtration preserved by multiplication into exact fixed-type restrictions.
The occupation estimates control both the total rational multiplicity of
the relations and the logarithmic cost of their auxiliaries. The circulation rounding preserves
integer resource balances, while its entropy estimate remains valid for
arbitrarily small or zero occupations. The separator applies to any
matched sectors sharing a leg. Each use in another tensor argument would
still need its own compatibility and cost analysis.

The threefold balance and pooled recursion have precise, limited effects.
Alternating row orientations make three cyclic copies sufficient and reduce
the dimensions at fixed parameters. This alone gives no improvement in
the exponent: threefold power \(2k\) reproduces the charged scalar counts
of sixfold power \(k\) for the same types and reuse count. Pooling the
recursive subproblems removes the ceiling in the branch ratio, replacing
\(\lceil r/m\rceil\) by \(r/m\) in the exponent. It does not remove
interpolation evaluations, catalyst supply, or linear transformations.

The scale bound applies to the relative source dimension \(A\). It does
not bound the dimension of the final rational seed, the size of its
coefficients, or the time needed to write them out. The seed procedure
terminates because it uses a fixed coefficient field and a strict gap
between the constructed exponent and the requested allowance. No new
small-rank identity is claimed. Every finite member of the prescribed
family has exponent strictly above \(9/4\), and the proof supplies no
fixed recursion with cost \(O(n^{9/4})\).

The allowance and coefficient set are fixed before the input dimension
grows. Although Theorem~\ref{thm:quantitative} gives a rate in the profile
index, it is not a construction-time or bit-complexity bound. A uniform
algorithm that changes its allowance with the input dimension would require
conventions for generating and loading its rational constants and a bound
on that setup cost. Exact arithmetic also gives no numerical stability
guarantee. The finite examples do not establish a useful matrix-size range.

Further reductions in finite cost can be assessed through the separate
terms of the exponent formula. A smaller separation degree would reduce
interpolation; a cheaper decomposition of the retained factors would reduce
catalyst supply; and different transition rules might reduce the occupation
and type costs. Such changes would need explicit maps or proved estimates.
The present analysis supplies no general lower bound that excludes them.

\subsection*{Supporting calculations and source files}

The proof is contained in the formulas and arguments in this paper.
The accompanying supplement, \path|finite-extraction-supplement-v2.zip|
(version 2, 7 October 2026), includes exact-arithmetic scripts, component
checks, and the saved interval certificates used for the finite examples.
The determinant and cyclic maps are checked in
\path|experiments/multiway_determinant.py| and
\path|experiments/cyclic_profile_extraction.py|. The occupation and rounded
counts are generated by \path|experiments/multiway_profile_flows.py| and
\path|experiments/multiway_rounded_flow.py|. The finite-cost evaluator is
\path|experiments/multiway_finite_costs.py|, with its saved report in
\path|results/multiway-finite-costs.json|. The supplement also retains the
binary-certificate calculations underlying the appendix comparisons.

These files are supporting checks, not assumptions of the main theorem.
The composed coefficient maps are specified mathematically; their full
arrays are not stored. The reported finite bounds come from saved rational
interval calculations, not from timing or executing the represented large
matrix algorithms. No numerical benchmark is used to infer an arithmetic
exponent.
